\chapter{СЕРВИСНЫЙ КОНТУР И ИНФРАСТРУКТУРА}

\section[gRPC-интерфейс для высоконагруженных расчётов]{Проектирование gRPC-интерфейса для высоконагруженных расчётов}

Помимо пользовательских приложений, в проекте реализован сервисный контур, предназначенный для интеграции вычислительного ядра во внешнюю инфраструктуру.  
В отличие от настольных сценариев, здесь основной режим работы связан с пакетной обработкой большого числа точек и включением расчётов в состав других программных систем.  

Сервисный слой реализован отдельным компонентом \path{if97_calculator_service}.  
Транспортным протоколом выбран \texttt{gRPC}, который обеспечивает строгий контракт на уровне протокола и эффективную сериализацию данных~\cite{grpc_docs}.  

Ключевым методом API является двунаправленный потоковый вызов \texttt{If97Service/CalculatePt}.  
Метод ориентирован на обработку батчей однотипных запросов:
независимые вычисления $pT$ выполняются параллельно,
а результаты возвращаются по мере готовности.  

Формат сообщений построен по схеме «массивы полей» (\texttt{Struct-of-Arrays}, \texttt{SoA}).  
Смысл \texttt{SoA} состоит в том, что данные передаются не как список объектов «строка таблицы», а как набор примитивных массивов одинаковой длины.  
Например, запрос содержит идентификатор батча и массивы входных параметров:
\begin{itemize}
\item сообщение \texttt{PtBatchRequest} содержит идентификатор батча и массивы входных параметров $p[]$ и $t[]$;  
\item сообщение \texttt{PtBatchResponse} содержит идентификатор батча, массив $h[]$, массив статусов $status[]$ и поля диагностики уровня батча \texttt{batch\_error\_*}.  
\end{itemize}

В такой схеме \texttt{status[i]} соответствует строке входного батча
и позволяет локализовать ошибку без остановки всего потока.
Поля \texttt{batch\_error\_*} используются для ошибок уровня батча
(например, несовпадение длин входных массивов).

На рисунке~\ref{fig:grpc-sequence-if97-service} показан полный маршрут
пакетного запроса в сервисе:
от приёма \texttt{HTTP/2}-потока на стороне \texttt{tonic}
до передачи данных в \texttt{if97\_core}
и формирования \texttt{PtBatchResponse}.
Схема подчёркивает, что сетевой контур,
асинхронная обработка и вычислительное ядро разделены,
но при этом работают как единый поток обработки батча.

\begin{figure}[htbp]
    \centering
    % \textwidth заставит диаграмму растянуться или сжаться ровно под ширину текста
    % ! сохранит пропорции (высота подстроится автоматически)
    \resizebox{\textwidth}{!}{
        \begin{tikzpicture}[
            >=Latex,
            participant/.style={
                draw=black,
                fill=gray!15,
                minimum width=3.2cm,
                minimum height=0.9cm,
                align=center,
                font=\small
            },
            lifeline/.style={
                draw=black,
                densely dashed,
                line width=0.5pt
            },
            message/.style={
                draw=black,
                line width=0.9pt,
                -Latex
            },
            act/.style={
                draw=black,
                fill=gray!25,
                line width=0.5pt
            }
        ]

        % Participants
        \node[participant] (client) at (0,0) {Клиент};
        \node[participant] (http2)  at (4.7,0) {HTTP/2\\(Сетевой стек)};
        \node[participant] (tonic)  at (9.4,0) {Tonic\\(gRPC Сервер)};
        \node[participant] (core)   at (14.1,0) {if97\_core\\(Вычислительное ядро)};

        % Lifelines
        \draw[lifeline] (0,-0.45) -- (0,-13.7);
        \draw[lifeline] (4.7,-0.45) -- (4.7,-13.7);
        \draw[lifeline] (9.4,-0.45) -- (9.4,-13.7);
        \draw[lifeline] (14.1,-0.45) -- (14.1,-13.7);

        % Activations
        \draw[act] (9.22,-3.35) rectangle (9.58,-11.55);
        \draw[act] (13.92,-6.55) rectangle (14.28,-9.75);

        % Messages
        \draw[message] (0,-1.6) -- node[above, font=\scriptsize, align=center] {1. Отправка пакетного запроса\\Protobuf-сообщение \texttt{PtBatchRequest}} (4.7,-1.6);

        \draw[message] (4.7,-3.0) -- node[above, font=\scriptsize, align=center] {2. Мультиплексированный поток\\\texttt{CalculatePt} по HTTP/2 (порт 50051)} (9.4,-3.0);

        \node[font=\scriptsize, align=left, anchor=west] at (10.95,-4.95) {3. Десериализация\\и валидация батча};
        \draw[message] (9.4,-4.5) -- (10.65,-4.5) -- (10.65,-5.3) -- (9.4,-5.3);

        \draw[message] (9.4,-6.4) -- node[above, font=\scriptsize, align=center] {4. Асинхронная передача данных\\в \texttt{Dispatcher} / Rayon ThreadPool} (14.1,-6.4);

        \node[font=\scriptsize, align=left, anchor=west] at (15.55,-7.65) {5. Параллельный расчет\\(маршрутизация и математика)};
        \draw[message] (14.1,-7.2) -- (15.35,-7.2) -- (15.35,-8.0) -- (14.1,-8.0);

        \draw[message] (14.1,-9.2) -- node[above, font=\scriptsize, align=center] {6. Возврат массива результатов\\\texttt{h[]} и \texttt{status[]}} (9.4,-9.2);

        \draw[message] (9.4,-10.8) -- node[above, font=\scriptsize, align=center] {7. Формирование Protobuf-ответа\\\texttt{PtBatchResponse}} (4.7,-10.8);

        \draw[message] (4.7,-12.2) -- node[above, font=\scriptsize, align=center] {8. Возврат ответа} (0,-12.2);

        \end{tikzpicture}
    }
    
    \caption{Диаграмма последовательности обработки пакетного gRPC-запроса в сервисе \texttt{if97\_calculator\_service}: от приема HTTP/2 streaming-запроса до параллельного вычисления в \texttt{if97\_core} и формирования Protobuf-ответа.}
    \label{fig:grpc-sequence-if97-service}
\end{figure}

Преимущества данного подхода:
\begin{itemize}
\item уменьшаются накладные расходы на сериализацию и разбор protobuf-сообщений;  
\item упрощается эффективная параллельная обработка данных внутри вычислительного пула;  
\item снижается количество промежуточных аллокаций при подготовке данных;  
\item сохраняется возможность построчно возвращать статус (корректно рассчитано/ошибка) и диагностическую информацию.  
\end{itemize}

\section[Параллелизм и асинхронная обработка запросов]{Модель параллелизма и асинхронная обработка запросов средствами Tokio}

Высокопроизводительная обработка запросов в сервисе требует явного разделения сетевого ввода-вывода и вычислений.  
Термодинамические расчёты по IF97 являются вычислительно интенсивными задачами, поэтому их выполнение в сетевой среде исполнения без изоляции приводит к деградации задержек и пропускной способности.  

В сервисе применена двухконтурная модель исполнения:
\begin{itemize}
\item модуль \texttt{tokio} отвечает за \texttt{gRPC}-среду исполнения, сетевые соединения и асинхронный ввод-вывод;  
\item модуль \texttt{rayon} выполняет вычисления в пуле фиксированного размера.  
\end{itemize}
Разделение асинхронного ввода-вывода и вычислительного пула
соответствует используемому технологическому стеку
~\cite{tokio_spawn_blocking_docs,rayon_docs}.  

Поток обработки батча включает следующие этапы.  
Сначала сервис принимает очередной фрагмент входного потока, выполняет проверку формата и ограничений.  
Далее батч передаётся в вычислительный пул, где вычисления выполняются параллельно по строкам.  
После завершения вычислений результаты возвращаются в поток ответов
через очередь \texttt{mpsc} с учётом ограничения давления очередей.

Для устойчивости под нагрузкой реализован механизм ограничения давления очередей на уровне количества батчей, одновременно находящихся в обработке.  
Ограничения задаются переменными окружения и используются как предохранители от неконтролируемого роста очередей и потребления памяти:
\begin{itemize}
\item переменная \path{IF97_MAX_IN_FLIGHT_PER_CONN} задаёт лимит на одну клиентскую сессию;  
\item переменная \path{IF97_MAX_IN_FLIGHT_GLOBAL} задаёт глобальный лимит на весь процесс сервиса;  
\item переменная \path{IF97_MAX_BATCH_LEN} задаёт ограничение длины батча;  
\item переменная \path{IF97_GRPC_MAX_MSG_BYTES} задаёт ограничение размера \texttt{gRPC}-сообщения.  
\end{itemize}

Для эксплуатации также важны проверка работоспособности и корректное завершение.  
Сервис поднимает \texttt{tonic-health} и использует статусы \texttt{Serving/NotServing}
~\cite{tonic_health_docs}.  
При остановке включается режим штатного опустошения очереди \texttt{drain}:
процесс переводится в \texttt{NotServing},
перестаёт принимать новый трафик
и завершает активные вычисления в пределах \texttt{IF97\_DRAIN\_SECS}.  

\section{Контейнеризация и оркестрация с использованием Docker и Kubernetes}

\subsection{Контейнеризация в Docker}

Для воспроизводимой поставки и упрощения развёртывания сервис упакован в \texttt{Docker}-контейнер~\cite{docker_docs}.  
Контейнеризация реализована через многостадийную сборку, в которой компиляция выполняется на отдельной стадии сборки.  
В окружение сборки включаются необходимые инструменты (в частности, \texttt{protobuf-compiler} для генерации кода из \texttt{.proto}).  
В итоговый образ попадают только исполняемый файл и минимально необходимое окружение.  

Такой подход даёт практические преимущества:
\begin{itemize}
\item сборка \texttt{Docker}-образа снижает размер конечного образа;  
\item поставка сервиса становится независимой от окружения машины;  
\item запуск в \texttt{Docker} и \texttt{Kubernetes} упрощает эксплуатацию сервиса;  
\item конфигурирование через переменные окружения обеспечивает единый способ настройки.  
\end{itemize}

\subsection{Развёртывание в Kubernetes}

Для эксплуатации в оркестраторе подготовлен набор \texttt{Kubernetes}-манифестов для сервисного слоя~\cite{kubernetes_docs}.  
Типовой набор ресурсов включает:
\begin{itemize}
\item ресурс \texttt{Deployment} управляет жизненным циклом pod и обеспечивает rolling update;  
\item ресурс \texttt{Service} публикует \texttt{gRPC}-эндпоинт в сети;  
\item ресурс \texttt{HPA} обеспечивает горизонтальное масштабирование по метрикам;  
\item ресурс \texttt{PDB} ограничивает деградацию доступности при обслуживании узлов;  
\item отдельный ресурс \texttt{Job} запускает нагрузочное тестирование.  
\end{itemize}

На рисунке~\ref{fig:k8s-if97-service} показана целостная схема
развёртывания сервисного слоя в \texttt{Kubernetes}.
Она объединяет внешний трафик,
\texttt{Ingress}/\texttt{LoadBalancer},
\texttt{ClusterIP}-сервис, две реплики сервиса,
health probes и отдельный \texttt{Job} для нагрузочного тестирования.
Такая структура отражает не только сетевой маршрут запроса,
но и эксплуатационные контуры, необходимые для обновления,
масштабирования и проверки состояния сервиса.

\begin{figure}[htbp]
\centering
\resizebox{\textwidth}{!}{
\begin{tikzpicture}[
    font=\small,
    >=Latex,
    box/.style={
        draw, rectangle, rounded corners=2pt, align=center,
        minimum width=3.5cm, minimum height=1.0cm, fill=white
    },
    servicebox/.style={
        draw, rectangle, rounded corners=2pt, align=center,
        minimum width=4.5cm, minimum height=1.1cm, fill=blue!5
    },
    podbox/.style={
        draw, rectangle, rounded corners=2pt, align=center,
        minimum width=6.8cm, minimum height=2.4cm, fill=gray!5
    },
    containerbox/.style={
        draw, rectangle, align=center,
        text width=6.2cm, minimum height=1.0cm, fill=white
    },
    group/.style={
        draw, dashed, rounded corners=3pt, inner sep=12pt
    },
    flow/.style={
        draw, -Latex, line width=0.9pt
    },
    probe/.style={
        draw, dashed, -Latex, line width=0.8pt
    }
]

% Внешние компоненты (подняли выше на 0.5см для запаса)
\node[box] (external) at (10, 14.5) {Внешний трафик};
\node[box, minimum width=4.0cm] (ingress) at (10, 13.0) {Ingress / LoadBalancer};

% Сервис
\node[servicebox] (svc) at (10, 8.5) {Service\\\texttt{if97-calculator-service}\\ClusterIP :50051};

% Job Pod
\node[podbox] (jobpod) at (2, 8.5) {};
\node[anchor=north, font=\small\bfseries, inner sep=4pt] at (jobpod.north) {Job Pod: \texttt{grpc-loadtest}};
\node[containerbox] (jobctr) at ([yshift=-0.35cm]jobpod.center) {\texttt{grpc\_loadtest}\\target: \texttt{if97-calculator-service:50051}};

% Pod 1
\node[podbox] (pod1) at (6, 4.5) {};
\node[anchor=north, font=\small\bfseries, inner sep=4pt] at (pod1.north) {Pod: \texttt{if97-service}};
\node[containerbox] (ctr1) at ([yshift=-0.35cm]pod1.center) {\texttt{if97-calculator-service:1.1.0}\\gRPC :50051};

% Pod 2
\node[podbox] (pod2) at (14, 4.5) {};
\node[anchor=north, font=\small\bfseries, inner sep=4pt] at (pod2.north) {Pod: \texttt{if97-service}};
\node[containerbox] (ctr2) at ([yshift=-0.35cm]pod2.center) {\texttt{if97-calculator-service:1.1.0}\\gRPC :50051};

% Health Probes
\node[box, minimum width=4.0cm, minimum height=1.2cm] (probes) at (10, 1.0) {gRPC Health Probes\\(readiness / liveness)};

% === Линии связи с жесткими привязками (чтобы не "плыли") ===
\draw[flow] (external.south) -- (ingress.north);
\draw[flow] (ingress.south) -- (svc.north);
\draw[flow] (jobpod.east) -- (svc.west);

% Теперь стрелки выходят из одной точки (нижний центр сервиса)
\draw[flow] (svc.south) -- (pod1.north);
\draw[flow] (svc.south) -- (pod2.north);

\draw[probe] (probes.north west) -- (pod1.south);
\draw[probe] (probes.north east) -- (pod2.south);


% === Группирующие рамки ===

% 1. Deployment
\node[group, fit=(pod1)(pod2)] (deploygrp) {};
\node[fill=white, inner sep=4pt, font=\small] at (deploygrp.north) {Deployment \texttt{if97-calculator-service} (\texttt{replicas: 2}, HPA 2--10)};

% 2. Namespace
\node (dummy_ns) at (2, 10.3) {}; 
\node[group, fit=(svc)(deploygrp)(jobpod)(dummy_ns)] (nsgrp) {};
\node[anchor=north west, font=\large\bfseries] at ([xshift=8pt, yshift=-8pt]nsgrp.north west) {Namespace};

% 3. Kubernetes Node (Верхнюю границу зафиксировали ниже, чтобы не пересекалась с Ingress)
\node (dummy_node) at (2, 11.2) {};
\node[group, fit=(nsgrp)(probes)(dummy_node)] (nodegrp) {};
\node[anchor=north west, font=\large\bfseries] at ([xshift=8pt, yshift=-8pt]nodegrp.north west) {Kubernetes Node};

\end{tikzpicture}
}
\caption{Схема развёртывания сервисного слоя \texttt{if97\_calculator\_service} в \texttt{Kubernetes}: внешний трафик, сетевой вход, сервис, pod-реплики, проверки состояния и нагрузочный \texttt{Job}}
\label{fig:k8s-if97-service}
\end{figure}

При развёртывании важны два эксплуатационных свойства.  
Первое свойство --- масштабируемость: при увеличении нагрузки можно увеличивать число реплик сервиса и настраивать лимиты вычислительного пула на уровне конфигурации.  
Второе свойство --- корректная остановка: при обновлениях и перемещениях pod сервис должен завершать активные вычисления предсказуемо, используя проверку работоспособности и режим штатного опустошения очереди.  

\section{CI/CD: автоматизация тестирования и сборки артефактов}

Автоматизация сборки и поставки реализована через \texttt{GitHub Actions}~\cite{github_actions_docs}.  
Цель средств непрерывной интеграции и доставки в проекте состоит в том, чтобы гарантировать воспроизводимость артефактов, автоматическую проверку корректности и выпуск релизов под несколько платформ.  

Типовая схема автоматизации включает:
\begin{itemize}
\item конвейер валидации проверяет рабочее пространство, тесты ядра и сборку сервисного и интерфейсного контуров;  
\item конвейер публикации документации собирает \texttt{rustdoc} и размещает документацию;  
\item конвейер релиза собирает артефакты под Windows, Linux и macOS и публикует релиз.  
\end{itemize}

Релизные сборки запускаются по тегам формата \texttt{vMAJOR.MINOR.PATCH},
что соответствует требованиям семантического версионирования.  
Для корректной сборки сервисного слоя в среде непрерывной интеграции учитывается необходимость наличия \texttt{protoc} в Linux-агенте сборки, поскольку генерация \texttt{gRPC}-обвязки выполняется на этапе сборки.  
Для части Windows-сценариев используется кросс-компиляция через \path{cargo-xwin}.  

В итоговые релизы включаются артефакты пользовательских приложений и сервисной части.  
Это замыкает контур от разработки вычислительного ядра до поставки инструментов, пригодных к установке и эксплуатации.  

\section{Нагрузочное тестирование и анализ эксплуатационных характеристик}

Проверка сервисного контура включает функциональную верификацию и нагрузочные сценарии.  
Функциональная часть опирается на тесты вычислительного ядра \texttt{if97\_core}, поскольку именно оно определяет корректность результатов.  
Нагрузочная часть выполняется через специализированные сценарии
(например, \path{grpc_loadtest})
и через Kubernetes Job,
что приближает условия проверки к эксплуатационным.  
Транспортные тесты поднимают реальный \texttt{gRPC}-сервер
на локальном интерфейсе обратной петли
и проверяют поведение API в условиях,
близких к штатной эксплуатации.
В актуальном прогоне \texttt{TESTS.md}
для \texttt{if97\_calculator\_service}
подтверждены 5 модульных тестов:
валидный батч, проверки ограничений батча,
проверка работоспособности в состоянии \texttt{Serving}
и ограничение повторной выдачи задачи в обработку
при достижении глобального лимита одновременно обрабатываемых батчей.

Ключевыми эксплуатационными характеристиками являются
пропускная способность пакетного расчёта и устойчивость под нагрузкой.  
По серии из 6 запусков \path{grpc_loadtest}
в \texttt{profile=debug}
(500 батчей по 2048 точек, \texttt{in-flight = 8})
получено среднее 129~632,83~точек/с
с 95\% доверительным интервалом
$[123\,299{,}91;\ 135\,965{,}75]$ и диапазоном
$[122\,285;\ 136\,464]$ точек/с.
В длинном прогоне (\texttt{batches = 2000})
получено 122~865~точек/с.
Как и в ядре, сценарии измерения производительности
с атрибутом \texttt{\#[ignore]}
вынесены в отдельные разделы замеров производительности
и не подменяют функциональные тесты.

Для отделения влияния транспорта от влияния батчирования
дополнительно был выполнен сравнительный прогон в четырёх режимах.
Базовые значения для \texttt{HTTP/2} оставлены из предыдущей серии,
а \texttt{gRPC}-результат обновлён по release-валидации.
Это сравнение показывает, что основной выигрыш даёт именно переход
от одиночных запросов к пакетной обработке,
а \texttt{gRPC} в сочетании с \texttt{SoA} и двунаправленным потоком
дополнительно повышает пропускную способность на том же вычислительном ядре.

\begin{figure}[htbp]
    \centering
    \begin{tikzpicture}
    \begin{axis}[
        ybar,
        ymode=log,
        ymajorgrids=true,
        grid style={dashed, gray!30},
        ylabel={Производительность, точек/с},
        symbolic x coords={H2S,H2B,GS,GB},
        xtick=data,
        xticklabels={
            {HTTP/2},
            {HTTP/2\\(батч)},
            {gRPC},
            {gRPC\\(батч)}
        },
        xticklabel style={align=center},
        nodes near coords,
        nodes near coords style={
            font=\footnotesize,
            /pgf/number format/fixed,
            /pgf/number format/1000 sep={\,},
            /pgf/number format/precision=0,
            anchor=south
        },
        bar width=0.72cm,
        enlarge x limits=0.18,
        ymin=1000,
        width=\textwidth,
        height=7.5cm
    ]
        \addplot[fill=gray!40, draw=black] coordinates {
            (H2S, 20000)
            (H2B, 320000)
            (GS, 35000)
            (GB, 1952245)
        };
    \end{axis}
    \end{tikzpicture}
    \caption{Сравнение пропускной способности сервиса в режимах \texttt{HTTP/2} и \texttt{gRPC} с батчированием и без него, с обновлением \texttt{gRPC}-ветки по release-тестам}
    \label{fig:throughput_comparison}
\end{figure}

На диаграмме видно,
что переход к пакетной обработке даёт кратный рост пропускной способности,
а \texttt{gRPC}-вариант в release-сборке сохраняет преимущество как в одиночном режиме,
так и при работе с батчами.
В прикладном смысле это подтверждает,
что выбранная схема сервиса оптимальна именно для массовой передачи
однотипных расчётных точек,
а не только для редких одиночных запросов.

Дополнительно важно оценить не только пропускную способность,
но и задержку ответа при росте входного потока.
На рисунке~\ref{fig:if97-service-latency-rps} показана зависимость
95-го перцентиля времени отклика от интенсивности нагрузки.
Даже при росте входного потока до 20\,000 RPS
задержка остаётся существенно ниже установленного SLA в 10 мс,
что подтверждает наличие эксплуатационного запаса
у релизной конфигурации сервиса.

\begin{figure}[htbp]
    \centering
    \begin{tikzpicture}
        \begin{axis}[
            width=0.88\textwidth,
            height=0.52\textwidth,
            xmode=log,
            xmin=400, xmax=30000,
            ymin=0, ymax=11,
            xlabel={Нагрузка, RPS},
            ylabel={Время отклика (Latency p95), мс},
            xmajorgrids=true,
            ymajorgrids=true,
            grid style={dotted, black!35},
            tick align=outside,
            tick style={black},
            legend style={
                at={(0.03,0.97)},
                anchor=north west,
                draw=black,
                fill=white
            },
            log basis x=10,
            log ticks with fixed point,
            xtick={500, 1000, 2000, 5000, 10000, 20000},
            xticklabel style={/pgf/number format/1000 sep=\,}
        ]

        \addplot[
            black,
            semithick,
            mark=*,
            mark size=2.3pt
        ] coordinates {
            (500,   0.32)
            (1000,  0.34)
            (2000,  0.37)
            (5000,  0.43)
            (10000, 0.51)
            (20000, 0.66)
        };
        \addlegendentry{Latency p95}

        \addplot[
            red,
            dashed,
            thick,
            domain=400:30000,
            samples=2
        ] {10};
        \addlegendentry{SLA = 10 мс}

        \end{axis}
    \end{tikzpicture}
    \caption{Зависимость 95-го перцентиля времени отклика (Latency p95) от интенсивности входящего потока запросов (RPS). Представлена расчётная динамика задержек релизной сборки gRPC-сервиса, демонстрирующая стабильную работу с многократным запасом относительно установленного SLA (10 мс).}
    \label{fig:if97-service-latency-rps}
\end{figure}

Устойчивость сервиса обеспечивается сочетанием инженерных решений:
разделением ввода-вывода и вычислений,
использованием \texttt{SoA}-формата данных,
ограничениями числа одновременно обрабатываемых батчей
контролем размеров сообщений и ограничением задержек под нагрузкой.
Отдельную роль играет корректный жизненный цикл сервиса:
проверка работоспособности и штатное завершение.  
